%% LyX 2.4.4 created this file.  For more info, see https://www.lyx.org/.
%% Do not edit unless you really know what you are doing.
\documentclass[english]{beamer}
\usepackage{mathptmx}
\usepackage[T1]{fontenc}
\usepackage[latin9]{inputenc}
\usepackage{units}
\usepackage{amssymb}
\usepackage{graphicx}

\makeatletter
%%%%%%%%%%%%%%%%%%%%%%%%%%%%%% Textclass specific LaTeX commands.
% this default might be overridden by plain title style
\newcommand\makebeamertitle{\frame{\maketitle}}%
% (ERT) argument for the TOC
\AtBeginDocument{%
  \let\origtableofcontents=\tableofcontents
  \def\tableofcontents{\@ifnextchar[{\origtableofcontents}{\gobbletableofcontents}}
  \def\gobbletableofcontents#1{\origtableofcontents}
}

%%%%%%%%%%%%%%%%%%%%%%%%%%%%%% User specified LaTeX commands.
\usetheme{Warsaw}
% or ...

\setbeamercovered{transparent}
% or whatever (possibly just delete it)
\setbeamercovered{invisible} %makes overlays invisible


%gets rid of navigation symbols
\setbeamertemplate{navigation symbols}{}

%\usepackage{pgfpages}
%\pgfpagesuselayout{4 on 1}[a4paper, landscape, border shrink=7mm]
%\pgfpageslogicalpageoptions{1}{border code=\pgfusepath{stroke}}
%\pgfpageslogicalpageoptions{2}{border code=\pgfusepath{stroke}}
%\pgfpageslogicalpageoptions{3}{border code=\pgfusepath{stroke}}
%\pgfpageslogicalpageoptions{4}{border code=\pgfusepath{stroke}}

\makeatother

\usepackage{babel}
\begin{document}
\title[Approximating Functions with Polynomials]{Ch. 10.1 Applications of Power Series}
\subtitle{Calculus III | MTH 331~~|~~Section A~~|~~Fall 2012}
\author{Dr.\,James Ashe}
\titlegraphic{\includegraphics[height=1.5cm]{../../jcsu_bull.jpg}}
\institute{Johnson C. Smith University}
\date{{}}

\makebeamertitle

%\beamerdefaultoverlayspecification{<+->}
\begin{frame}{Power Series}

\begin{block}{Definition}

A \emph{power series }is an infinite series of the form 
\[
{\displaystyle \sum_{k=0}^{\infty}c_{k}x^{k}=c_{0}+c_{1}x+c_{2}x^{2}+\dots+c_{n}x^{n}+c_{n+1}x^{n+1}\dots}
\]

or, more generally, 
\[
{\displaystyle \sum_{k=0}^{\infty}c_{k}(x-a)^{k}=c_{0}+c_{1}(x-a)+\dots+c_{n}(x-a)^{n}+c_{n+1}(x-a)^{n+1}\dots}
\]

where each $c_{k}$ and the \textbf{center }of the series $a$ are
constants.
\end{block}
\end{frame}

\begin{frame}{}

We can estimate a function, say $f(x)=\ln(x+1)$ at $x=0$ using its
derivative, $f'(x)=\frac{1}{x+1}$, $f''(x)=-\frac{1}{(x+1)^{2}}$,
and $f'''(x)=\frac{2}{(x+1)^{3}}$. 
\begin{columns}


\column{6.5cm}
\begin{itemize}
\item []<2->\textcolor{blue}{$p_{1}(x)=f(a)+f(a)(x-a)$}
\item []<3->$p_{1}(x)=f(0)+f'(0)(x-0)$\uncover<4->{$\Rightarrow p_{1}(x)=0+\frac{1}{1}(x-0)=x$.}
\item []<5->\textcolor{blue}{$p_{2}(x)=p_{1}+\frac{1}{2}f''(a)(x-a)^{2}$}
\item []<6->$p_{2}(x)=(x)-\frac{1}{2}f''(0)(x-0)^{2}$\uncover<7->{$\Rightarrow p_{2}(x)=x-\frac{1}{2}x{}^{2}$}
\item []<8->$p_{3}(x)=x-\frac{1}{2}x^{2}+\frac{1}{3}x^{3}$
\end{itemize}

\column{5cm}

\includegraphics[width=5cm]{images/taylor_approx_ln(x+1)_0}<1-3|handout:0>\includegraphics[width=5cm]{images/taylor_approx_ln(x+1)_1}<4-6|handout:0>\includegraphics[width=5cm]{images/taylor_approx_ln(x+1)_2}<7|handout:0>\includegraphics[width=5cm]{images/taylor_approx_ln(x+1)_3}<8->
\end{columns}
\end{frame}

\begin{frame}{Taylor Polynomials}

In the last example $p_{3}$ was built on top of $p_{2}$. So the
following conditions are met:
\begin{itemize}
\item []<2->$p_{3}(a)=f(a),\,\,p_{3}'(a)=f'(a),\,\,p''(a)=f''(a),$\textrm{
and $\mbox{\ensuremath{p_{3}'''(a)=f'''(a)\mbox{.}}}$}
\end{itemize}
\uncover<3->{Solving for $c_{3}$ (similarly to how $c_{2}$ was
found), we get $c_{3}=\frac{f'''(a)}{3!}$. This gives the cubic polynomial
approximation,}
\begin{itemize}
\item []<4->$p_{3}(x)=f(a)+f'(a)(x-a)+\frac{f''(a)}{2!}(x-a)^{2}+\frac{f'''(a)}{3!}(x-a)^{3}\mbox{.}$
\end{itemize}
\begin{block}<5->{\textbf{Definition: }Taylor Polynomials}

Let $f$ be a function with $f',f'',\dots,f^{(n)}$ defined at $a$.
The $n^{\mbox{th}}$ \emph{order Taylor polynomial }for $f$ with
\textbf{center $a$ }is, $p_{n}(x)=f(a)+f'(a)(x-a)+\frac{f''(a)}{2!}(x-a)^{2}+\cdots+\frac{f^{(n)}(a)}{n!}(x-a)^{n}$.
\end{block}
\end{frame}

\begin{frame}{Approximations with Taylor Polynomials}

Ex. (see 10.1: 23) Use a Taylor polynomial of degree $3$ to approximate
$\frac{1}{\sqrt{1.08}}$.
\begin{itemize}
\item []<2->Use $f(x)=\frac{1}{\sqrt{1+x}}$ and $a=0$ since $\sqrt{1}$
is easy to find.
\item []<3->$p_{3}(x)=f(a)+f'(a)(x-a)+\frac{f''(a)}{2!}(x-a)^{2}$ and
$f'(x)=-\frac{1}{2}\frac{1}{(x+1)^{\nicefrac{3}{2}}}$, $f''(x)=\frac{3}{4}\frac{1}{(x+1)^{\nicefrac{5}{2}}}$.
\item []<4->So $p_{3}(x)=1-\frac{1}{2}\frac{1}{(1)^{\nicefrac{3}{2}}}(x)+\frac{3}{4}\frac{1}{(1)^{\nicefrac{5}{2}}}(x)$,
and \uncover<5->{$p_{3}(0.08)=1-\frac{1}{2}(0.08)+\frac{3}{4}(0.08)\approx1.02$}
\item []<5->Which is a decent approximation of $\sqrt{1.08}\approx1.03923$.
\end{itemize}
\end{frame}

\begin{frame}{}

\textbf{Homework }due 10/24

\textbf{10.1: }7-8, 13-16, 21-22, 27-28, 33-36, 41-42
\end{frame}

\end{document}
